\SourceLeaf{512}die noch mehr zu vermindern oder ganz wegzuschaffen sehr wünschenswerth wäre. Da Hr.\ Prof.\ \textsc{Bohnenbergers} zu diesem Zwecke angestellte Versuche, der Äusserung S.\ 392 zufolge, ohne Erfolg gewesen sind, und die Vermuthung zu begründen scheinen könnten, dass dies unmöglich sei, so hat mich dies zu einer besondern Untersuchung veranlasst, aus der sich, was mir sehr merkwürdig scheint, das Gegentheil ergeben hat.

Die vollkommne Wegschaffung der Farbenzerstreuung bei den Randstrahlen und den der Axe nächsten Strahlen ist nemlich allerdings möglich, oder bestimmter, es lässt sich ein Objectiv berechnen, welches alle Strahlen von zwei bestimmten Farben, sowohl diejenigen, welche in einer bestimmten Entfernung von der Axe, als die, welche unendlich nahe bei derselben (und zwar, wie hier immer vorausgesetzt wird, mit ihr parallel) auffallen, in Einen und demselben Punkt vereinigt. Dies Objectiv erhält eine von den bisher ausschliesslich angewandten ganz abweichende Form, so dass beide Linsen convex-concav werden und die convexen Flächen dem Gegenstande zukehren. Inzwischen obgleich hierdurch grössere Brechungen vorkommen als bei andern Einrichtungen, ist dennoch die übrig bleibende unvermeidliche Abweichung wegen der Gestalt noch sehr unbedeutend, und also die Vereinigung \textit{aller} auf das Objectiv parallel mit der Axe auffallenden Strahlen vollkommener als bei irgend einer andern Einrichtung. Es wäre daher wohl der Mühe werth, dass geschickte Künstler diese neue Form versuchten. Es kann vielleicht sein, dass gegenwärtig dabei noch \textit{practische} Schwierigkeiten statt finden; eine davon wird die sein, dass die Glasstücken, aus denen die Linsen geschliffen werden sollen, eine grössere Dicke haben müssen. Allein bei der immer fortschreitenden Vollkommenheit des technischen Theils der Dioptrik steht zu hoffen, dass Schwierigkeiten der Art zu besiegen sein werden, und dann ist es an der Mathematik, das Ideal der Form zur vollkommensten Vereinigung zu geben.

Die von mir geführte Rechnung soll übrigens bloss als Beispiel dienen, das Gesagte zu bestätigen, nicht aber dazu, dass Künstler diese Maasse genau befolgen sollen. Es ist unumgänglich nothwendig, dass für die Glasarten, aus denen ein vollkommenes Objectiv geschliffen werden soll, die Brechungs- und Zerstreuungsverhältnisse erst besonders mit möglichster Schärfe bestimmt, und die Maasse des Objectivs diesem gemäss von Neuem berechnet werden. In meiner Rechnung habe ich genau dieselben Zahlen zum Grunde gelegt, nach denen Hr.
\SourceLeaf{513}\expandafter\def\csname HistoricalNoteText513\endcsname{[Dicke der ersten Linse $=200$, der zweiten $=80$ Abstand zwischen beiden Linsen $=50$. Exponenten der Brechungsverhältnisse bezüglich für Kronglas und Flintglas]
\[
\left\{\begin{array}{lll}
1.525976&1.62173&\text{viol. Strahlen}\\
1.515162&1.60177&\text{mittl.}\quad\text{—}\\
1.504348&1.58181&\text{rothe}\quad\text{—}
\end{array}\right\}
\]
}Prof.\ \textsc{Bohnenberger} gerechnet hat: auch dieselbe Dicke und Entfernung der Linsen habe ich beibehalten\HistoricalNote{513}. Da aber bei der neuen Einrichtung die convexe Fläche der Flintglaslinse eine stärkere Krümmung hat, als die concave der Kronglaslinse, so können beide Linsen einander näher kommen (welches auch in einer andern hier nicht weiter auszuführenden Rücksicht vortheilhafter sein wird); ja, wenn die Künstler sonst keine Bedenklichkeit dagegen haben, kann der Zwischenraum ganz wegfallen, oder die Linsen können einander in der Axe berühren. Es versteht sich, dass dies einige Modification der Krümmungshalbmesser nach sich ziehen wird.

Es gehört nicht zu meiner Absicht, den mathematischen Theil dieser Untersuchung hier zu entwickeln. Ich bemerke nur, dass die Aufgabe, wenn man die Abweichung wegen der Gestalt nach \textsc{Eulers} Art betrachtet, und Dicke und Entfernung der Glaslinse bei Seite setzt, auf eine Gleichung des vierten Grades führt, welche zwei reelle Wurzeln hat. Die hieraus sich ergebende genäherte Auflösung dient zur Grundlage einer indirecten Rechnung, durch welche alles genau in Übereinstimmung gebracht wird. Für Mathematiker wird diese Andeutung hinreichen. Die eine reelle Wurzel jener Gleichung muss übrigens verworfen werden, weil mit ihr zu starke Krümmungen der Glasflächen zusammenhängen, und die unvollkommene Aufhebung wegen der Gestalt zu sehr fühlbar machen würden.

Das Resultat meiner Rechnung ist nun folgendes:

Wenn die Halbmesser der Reihe nach zu
\[
+3415{,}287;\quad -10133{,}007;\quad +4207{,}421;\quad -2807{,}320
\]
angenommen werden, so vereinigen sich die rothen und violetten Strahlen, sowohl die, welche unendlich nahe bei der Axe, als die, welche in der Entfernung $1083{,}687$ auffallen, alle in Einem Punkt der Axe, dessen Entfernung von der letzten Fläche $=28293{,}3$ wird. Wird jene Entfernung von der Axe, bei welcher der Einfallswinkel $18^0\,30'$ ist, als Halbmesser der Öffnung angenommen, so ist der Durchmesser der Öffnung sehr nahe $\tfrac{1}{13}$ der Brennweite. Um beurthei\SourceLeaf{514}len zu können, wie gross die noch übrig bleibende Abweichung wegen der Gestalt für die Strahlen zwischen dem Rande und der Axe wird, habe ich die Vereinigungsweiten für den Einfallswinkel $13^0$ berechnet und gefunden
\[
\begin{array}{ll}
28289{,}3&\text{für die rothen}\\
28290{,}0&\text{für die violetten Strahlen.}
\end{array}
\]

Ich kann nicht umhin, hier noch eine Erinnerung über eine Äusserung des Hrn.\ Prof.\ \textsc{Bohnenberger} in dem erwähnten Aufsatze beizufügen. Ich halte nemlich dafür, dass es am vortheilhaftesten ist, die Abweichung wegen der Gestalt genau für die Randstrahlen zu heben. Hr.\ Prof.\ \textsc{Bohnenberger} hat S.\ 279 dieses Verfahren wie mir deucht mit Unrecht getadelt. Man könnte, sagt er, wenn die Abweichung für die Randstrahlen genau gehoben sei, die Öffnung ohne Schaden der Deutlichkeit bis dahin vergrössern, wo die Abweichung wieder der grössten Abweichung der Zwischenstrahlen gleich werde, und es sei daher am vortheilhaftesten, die Abweichung nicht für die Randstrahlen, sondern für Strahlen zwischen dem Rande und der Axe zu heben. Dies würde allerdings wahr sein, wenn die übrigbleibende Abweichung jenseits und diesseits der Entfernung, für welche sie gehoben ist, \textit{einerlei Zeichen} hätte, was aber nicht der Fall ist. Man könnte zwar hiegegen mit einigem Schein einwenden, dass es bei der Längenabweichung auf das Zeichen gar nicht ankomme, und dass positive und negative Abweichungen eine und dieselbe Undeutlichkeit im Auge hervorbringen. Allein hiebei nähme man offenbar stillschweigend an, dass das Ocular immer genau für das deutliche Sehen desjenigen Bildes gestellt sei, welches durch die der Axe nächsten Strahlen hervorgebracht wird, und dies kann doch nicht eingeräumt werden. Man mag dies Bild immerhin das Hauptbild nennen: es fällt mit dem von den Randstrahlen hervorgebrachten Bilde zusammen, wenn die Abweichung für diese gehoben ist, und alle übrigen Bilder werden dann (wenigstens allgemein zu reden) jenseits oder diesseits des Hauptbildes liegen. Da man nun das Ocular immer so stellt, dass die Undeutlichkeit so klein wie möglich wird, so sieht. man gerade das Hauptbild am wenigsten deutlich, und jede Vergrösserung der Öffnung vergrössert auch die Undeutlichkeit. Eine ausführlichere Erörterung dieses Umstandes würde mich hier zu weit abführen.
\clearpage
\fancyhead[L]{\small BRIEF AN BRANDES.}
\SourceLeaf{515}\pdfbookmark[0]{Brief von Gauss an Brandes}{brief_brandes}
\expandafter\def\csname HistoricalNoteText515\endcsname{[Über die achromatischen Doppelobjective besonders in Rücksicht der vollkommenen Aufhebung der Farbenzerstreuung.]}
\begin{center}
\hrule\vspace{2mm}
\textsc{Gehler's} Physikalisches Wörterbuch. 1831. Artikel:\par
Linsenglas, Berechnungen über achromatische und aplanatische Linsengläser aus zwei Glaslinsen.
\vspace{2mm}\hrule\vspace{5mm}
{\large\itshape Brief von Gauss an Brandes.\par}
\end{center}
Auf Veranlassung Ihres Briefes habe ich eine freie Stunde auf den in jenem Aufsatze\HistoricalNote{515} am Ende kurz erwähnten Umstand gewandt. Der eigentliche Sinn der dortigen Bemerkung scheint nicht von allen ganz richtig aufgefasst zu sein, aber auch meine Angabe bedarf einer kleinen Modification. Ich finde nemlich jetzt durch eine \textit{tiefer eindringende} Untersuchung, dass die Undeutlichkeit, die in dem Ausdrucke für die Längen-Abweichung von der vierten Potenz des Abstandes der auffallenden Strahlen von der Axe abhängt, den möglich kleinsten Total-Einfluss hat, wenn man das Objectiv so construirt, dass diejenigen Strahlen, die unendlich nahe bei der Axe einfallen, und diejenigen, die in einer Entfernung $=R.\BareRoot{\tfrac65}$ auffallen würden, (wo $R=$ Radius des Objectivs ist) in \textit{einem} Punkte $A$ sich vereinigen, wobei das Ocular dann so steht, dass man denjenigen Punkt der Axe, wo die Strahlen, die in der Entfernung $=(\tfrac35-\tfrac{\BareRoot{6}}{10})R$ und $=(\tfrac35+\tfrac{\BareRoot{6}}{10})R$ von der Axe aufgefallen sind, sich alle vereinigen, deutlich sieht. Denken Sie Sich nemlich durch diesen Punkt eine auf die Axe senkrechte Ebene, so ist das Bild desto undeutlicher, je grösser der Kreis um $A$ ist, den die von einem Punkte des Objects auf das Objectivglas gefallenen Strahlen füllen, doch so, dass die Intensität der Strahlen an jeder Stelle dieses Kreises mit berücksichtigt werden muss. Hiebei ist nun einige Willkürlichkeit; ich halte für das zweckmässigste, hier nach denselben Principien zu verfahren, die der Methode der kleinsten Quadrate zum Grunde liegen. Ist nemlich $\mathrm{d}s$ ein Element dieses Kreises, $\rho$ die Entfernung des Elements von $A$, und $i$ die Intensität der Strahlen daselbst, so nehme ich an, dass $\int i\rho\rho\,\mathrm{d}s$ als das Maass der Total-Undeutlich\SourceLeaf{516}keit zu betrachten sei, und mache dies zu einem Minimum. Ich finde dabei folgende Resultate: 1. Construirte man das Objectiv so, dass dasjenige Glied der Längen-Abweichung, welches von dem Quadrate der Entfernung von der Axe abhängt, $=0$ wird, und setzte das Ocular so. dass $A$ dahin fällt, wo die der Axe unendlich nahen Strahlen diese schneiden, so sei der Werth dieses Integrals $=E$. 2. Stellte man aber \textit{bei derselben Einrichtung} das Ocular so, dass das Integral so klein wird, wie es bei dieser Einrichtung werden kann (wobei $A$ der Vereinigungspunkt der in der Entfernung $=R\BareRoot[4]{\tfrac12}$ auffallenden Strahlen sein wird), so ist das Integral $=\tfrac14E$. 3. Dagegen ist bei der obigen Einrichtung und der vortheilhaftesten Stellung des Oculars das Integral $=\tfrac{1}{100}E$, als absolutes Minimum. Obiges Resultat, dass nemlich mit dem Vereinigungspunkte der der Axe unendlich nahen Strahlen ein bloss \textit{fingirtes} Bild (von Strahlen aus grösserer Distanz von der Axe als der Halbmesser des Objectivs) vereinigt werden soll, ist anfangs sehr überraschend und paradox scheinend; aber bei näherer Betrachtung sieht man den eigentlichen Grund leicht ein. Jenes erste \textit{sogenannte} Hauptbild (von Strahlen sehr nahe bei der Axe) ist nemlich dabei gleichsam das Unwichtigste wegen seiner geringen Intensität, viel wichtiger ist, dass die Strahlen von den der Peripherie näheren Ringen des Objectivs \textit{unter sich} besser zusammen gehalten werden, was bei jener Einrichtung am besten erreicht wird. Es thut mir leid, dass die Grenzen eines Briefes jetzt grössere Ausführlichkeit nicht gestatten; der scharfe Calcül lässt sich nichts abstreiten und bei einem vagen Raisonnement übersieht man leicht einen wesentlichen Umstand; allein für den Kenner werden diese Winke schon zureichen.

Allgemein finde ich, dass immer \textit{bei der vortheilhaftesten Stellung des Oculars} jenes Integral $=\tfrac14E(1-\tfrac85\mu\mu+\tfrac23\mu^4)$ wird, wenn das Objectiv so construirt ist, dass Strahlen aus der Entfernung $\mu R$ von der Axe sich mit dem (oben sogenannten) Hauptbilde in einem Punkte vereinigen. Dies ist ein Minimum für $\mu=\BareRoot{\tfrac65}$ und ist dann $=\tfrac{1}{100}E$; für $\mu=1$ wäre es nur $=\tfrac{1}{60}E$ und für $\mu=$ unendlich klein, $=\tfrac14E$. Nicht allein hat also hienach \textsc{Bohnenberger} Unrecht, sondern auch ich habe damals Unrecht gehabt, aber insofern, als ich noch nicht weit genug von \textsc{Bohnenberger} abgewichen bin. Ich hatte damals bloss die \textit{ganze Grösse} des undeutlichen Bildes berücksichtigt, ohne auf die ungleiche Intensität der einzelnen Theile Rücksicht zu nehmen.

\clearpage
\fancyhead[L]{\small BERICHTIGUNG DER SCHNEIDEN EINER WAAGE.}
\SourceLeaf{517}\pdfbookmark[0]{Berichtigung der Schneiden einer Waage}{schneiden_waage}
\begin{center}
{\Large [BERICHTIGUNG DER SCHNEIDEN EINER WAAGE.]\par}
\vspace{7mm}\hrule\vspace{2mm}
Göttingische gelehrte Anzeigen. 1837 März 13.
\vspace{2mm}\hrule\vspace{5mm}
\end{center}
In der Sitzung der Königl.\ Gesellschaft der Wissenschaften vom 28. Januar nahm der Hofr.\ \textsc{Gauss} von der Vorlesung des Hrn.\ Prof.\ \textsc{Weber}, über welche im 22. Stücke dieser Blätter Bericht abgestattet ist, Veranlassung, einen Vortrag über einen nahe verwandten Gegenstand zu halten, von welchem wir den Hauptinhalt hier zur Anzeige bringen.

Er betrifft eine neue Berichtigungsmethode zur Erfüllung einer wesent\-lichen Bedingung bei den feineren Hebelwaagen, deren Wichtigkeit bisher nicht genug gewürdigt zu sein scheint. Solche Waagen haben drei prismatische Schneiden; die eine nach unten gekehrte, in der Mitte des Waagebalkens, ruht auf einem harten horizontalen Lager von Stein oder Stahl, und dient als Drehungsaxe bei dem Spiel des Waagebalkens; die beiden andern an den Enden des Waagebalkens sind aufwärts gerichtet, und auf jeder derselben schwebt das Tragestück, woran die Waageschale hängt. Die Tragestücke selbst sind von gehärtetem Stahl, und ihre unteren, auf den Schneiden aufliegenden Flächen vollkommen plan und hochpolirt.

Eine wesentliche Bedingung ist nun, dass diese beiden äussern Schneiden mit der mittleren parallel sein sollen. In der That. da vor jedem Umtausch der Gewichte in einer Schale die Waage erst gehemmt und dabei das Tragestück von der Schneide abgehoben wird, so ist nie darauf zu rechnen, dass sich nach Aufhebung der Hemmung das Tragestück \textit{genau} wieder eben so auf die Schneide legt, wie zuvor: dies ist zwar unschädlich, wenn die betreffende Schneide mit der mittleren parallel ist, verursacht aber ein verändertes Moment, wenn eine Divergenz der Schneiden statt findet. Eine unvollkommene Berichtigung in dieser Beziehung ist eine Hauptursache, warum bei oft wiederholten Wägungen zuweilen bedeutend grössere Abweichungen in den Resultaten sich zeigen, als man
\SourceLeaf{518}sonst von der vortrefflichen Arbeit und der Empfindlichkeit einer Waage erwarten sollte.

Die Mittel, deren sich die Künstler zur Berichtigung des Parallelismus der Schneiden bisher gewöhnlich bedient haben, sind nicht geeignet, alle zu wünschende Schärfe zu geben; auch ist es, bei feinen Waagen wie bei astronomischen Instrumenten, nicht der Verfertiger, von dem man die feinste Berichtigung zu fordern hat, sondern diese kommt dem zu, der die Waage gebraucht.

Das Verfahren, dessen sich der Hofr.\ \textsc{Gauss} zu dieser Berichtigung mit dem besten Erfolge bedient hat, beruht auf folgenden Gründen.

Bei den Schwingungen des Waagebalkens verändert die zu prüfende äussere Schneide zwar ihre Lage im Raume; diese verschiedenen Lagen sind aber alle unter einander parallel, wenn diese Schneide mit der (ruhenden) mittleren parallel ist. Anders verhält es sich dagegen, wenn die äussere Schneide der mittleren nicht parallel ist. Nehmen wir, um die Vorstellung zu fixiren, an, dass die äussere Schneide zwar mit der mittleren in Einer Ebene liege, dass aber die Richtungen der beiden Schneiden abwärts vom Beobachter divergiren. In diesem Falle wird bei dem Spiele des Waagebalkens die äussere Schneide sich auf einer Kegelfläche bewegen; ihr abwärts gekehrtes Ende wird, relativ gegen das nähere Ende, steigen oder sinken, so wie der Hebelarm, an welchem diese Schneide sich befindet, steigt oder sinkt. Dasselbe wird von dem die Schneide stets berührenden Tragestücke gelten.

Welcher von beiden Fällen nun statt finde, lässt sich erkennen, wenn auf dem Tragestücke ein Planspiegel befestigt ist. Am vortheilhaftesten ist es, diesen Spiegel so anzubringen, dass seine Ebene nahe senkrecht zu der Schneide ist, obwohl man darin nicht zu ängstlich zu sein braucht. In dem ersten der beiden Fälle bleibt der Spiegel, während des Spiels des Waagebalkens, sich selbst parallel, im zweiten nicht; im ersten Falle wird also das Bild eines in schicklicher Entfernung vor dem Spiegel sich befindenden Gegenstandes unverrückt bleiben, im zweiten hingegen (wie man leicht übersieht), mit dem betreffenden Hebelarme steigen oder sinken. Das umgekehrte würde statt finden, wenn die beiden Schneiden anstatt abwärts vom Beobachter zu divergiren, convergiren, es würde dann nemlich mit dem Steigen des Waagebalkenarmes ein Sinken des Bildes, und umgekehrt, verbunden sein.

Nun lässt sich, wenn der Spiegel ein sehr vollkommner ist, selbst eine
\SourceLeaf{519}äusserst kleine Verrückung des Bildes sicher und scharf mit einem Fernrohre erkennen. Der Hofr.\ \textsc{Gauss} gebrauchte als Gegenstand eine etwa 5 Meter vor dem Spiegel vertical aufgerichtete, in Millimeter eingetheilte Scale; das 35 mal vergrössernde Fernrohr stand in nahe eben so grosser Entfernung. Es erschien so das Bild eines Millimeters etwa 20 Secunden gross, wovon man noch Zehntel schätzen kann. So lange die Schneide noch nicht vollkommen berichtigt war, ging das Bild der Scale an dem Fadenkreuze des Fernrohrs auf das regelmässigste auf und ab, wie der Waagebalken seine Schwingungen machte.

Für mathematisch gebildete Leser bedarf es blos der Andeutung, dass auf diese Weise nicht blos erkannt werden kann, nach welcher Seite eine Divergenz statt findet, sondern auch, hinreichend genau, wie gross dieselbe ist, wodurch, verbunden mit der Kenntniss der Weite der Gewinde der Correctionsschrauben, das Correctionsgeschäft in einen sichern Gang gebracht wird.

Der Vollständigkeit wegen mögen noch ein Paar andere Umstände hier erwähnt werden.

Wenn man einen etwas grossen Spiegel anwendet (der vom Hofr.\ \textsc{Gauss} gebrauchte, auf das Tragestück vermittelst einer eigenen Vorrichtung befestigte, hat 75 Millimeter Höhe), so ist es nothwendig, die Schalen mit hinlänglich schweren Gewichten zu belasten, weil sonst das Tragestück seitwärts umschlagen würde.

Es ist oben vorausgesetzt, dass die zu prüfende äussere Schneide mit der mittleren in Einer Ebene liege, also, wenn man die mittlere genau horizontal gestellt hat, bei horizontalem Stande des Waagebalkens gleichfalls horizontal sei, und nur etwa seitwärts divergire. Gewöhnlich wird aber diese Voraussetzung auch nicht in äusserster Schärfe statt finden, sondern, die äussere Schneide bei jener Stellung etwas geneigt, oder das eine Ende etwas höher sein können als das andere. Man erkennt dies, bei der beschriebenen Prüfungsmethode, daran, wenn beim Steigen des Waagebalkenarmes das Spiegelbild sich zugleich seitwärts, und beim Sinken nach der entgegengesetzten Seite bewegt. Inzwischen muss bemerkt werden, dass dieser Fehler, wenn er vorhanden ist, an einer Waage von einem geschickten Künstler jedenfalls viel zu klein sein wird, um einen noch merklichen Fehler in den Resultaten der Wägungen hervor zu bringen, und dass man daher auch bei den besten Waagen keine Correctionsmittel zur Wegschaffung dieses Theils des Nicht-Parallelismus angebracht hat.

\clearpage
\fancyhead[L]{}
\SourceLeaf{521}\pdfbookmark[0]{Physikalische Beobachtungen}{physikalische_beobachtungen}
\thispagestyle{empty}
\begin{center}
\vspace*{44mm}
{\LARGE PHYSIKALISCHE\par}
\vspace{24mm}
{\LARGE BEOBACHTUNGEN.\par}
\end{center}
\clearpage
